\chapter{Случайни величини}

\section{Дефиниция и~основни свойства}

\begin{rembox}[Идея]
Случайната величина е, грубо казано, функция/изображение, което на
всяко елементарно събитие $\omega$ от $\Omega$ съпоставя някакво
реално число. Резултатът е случаен, защото изходът $\omega$ е
случаен, но самата функция е детерминистична - след като
експериментът се случи, $X$ приема една напълно определена числова
стойност.
\end{rembox}

\begin{defbox}[~(Случайна величина)]
Нека $V=(\Omega,\mathcal{A},\mathbb{P})$ е вероятностно пространство.
Функцията $X:\Omega\to\mathbb{R}$ е \textbf{случайна величина}
(стохастична), ако за всяко реално число $x\in\mathbb{R}$
праобразът на интервала $(-\infty,x)$ е събитие, т.е.
\[
  X^{-1}\big((-\infty,x)\big) = \{\omega\in\Omega : X(\omega)<x\} \in \mathcal{A}.
\]
Това свойство се нарича \textbf{измеримост} на $X$ спрямо
$\sigma$-алгебрата $\mathcal{A}$ - трябва да можем да кажем каква е
вероятността $X$ да е по-малко от $x$, за всяко~$x$.
\end{defbox}

\begin{rembox}[Забележка]
За да докажем, че $X$ е случайна величина, е достатъчно да проверим,
че $X^{-1}((-\infty,x))\in\mathcal A$ за всяко $x\in\mathbb R$. Оттам
насетне, за всяко борелово множество $B\in\mathcal B(\mathbb R)$
следва по дефиниция, че $X^{-1}(B)\in\mathcal A$ - например за
$B=[a,b]$, $B=\{a\}$, $B=(a,b)$ и~т.н.
\end{rembox}

\begin{thmbox}[~(Свойства на случайните величини)]{Теорема}
Нека $V$ е вероятностно пространство и~$X,Y:\Omega\to\mathbb R$ са
случайни величини. Тогава:
\begin{enumerate}[label=\arabic*)]
  \item $aX \pm bY$ е случайна величина за всички $a,b\in\mathbb R$;
  \item $cX$ е случайна величина за всяко $c\in\mathbb R$ (частен случай на 1) при $a=c$, $b=0$);
  \item $XY$ е случайна величина;
  \item $\dfrac{X}{Y}$ е случайна величина, ако $\mathbb P(Y_\omega{=}0):=\mathbb P(\{\omega\in\Omega:Y(\omega)=0\})=0$.
\end{enumerate}
\end{thmbox}

\begin{proofbox}
\textbf{1)} Нека първо $Z_1=aX$ при $a>0$. За $\{aX<y\}$, $y\in
\mathbb R$: $\{aX<y\} \stackrel{a>0}{=} \{X<\tfrac{y}{a}\}\in
\mathcal{A}$, понеже $X$ е случайна величина. Следователно $Z_1=aX$
е случайна величина (аналогично се разглеждат случаите $a<0$ и~$a=0$). По същия начин $Z_2=bY$ е случайна величина.

Остава да докажем, че $Z=Z_1+Z_2$ е случайна величина, т.е.\ че
$\{Z_1+Z_2<c\}\in\mathcal A$ за всяко $c\in\mathbb R$. Ключовото
твърдение е:
\[
  \{Z_1+Z_2<c\} = \bigcup_{q\in\mathbb Q} \Big(\{Z_1<q\}\cap\{Z_2<c-q\}\Big).
\]
(Обединението е по рационалните числа - изброимо, а~не по
$\mathbb{R}$ - точно защото $\sigma$-алгебрата гарантира затвореност
само относно \emph{изброимо} обединение.) За да докажем равенството
на множествата $A=\{Z_1+Z_2<c\}$ и~$B=\bigcup_q(\dots)$, показваме
двете влагания:
\begin{itemize}[nosep]
  \item $B\subseteq A$: ако $\omega\in B$, съществува $q$ такова, че
  $Z_1(\omega)<q$ и~$Z_2(\omega)<c-q$; събирайки - $Z_1(\omega)+
  Z_2(\omega) < q+(c-q)=c$, т.е.\ $\omega\in A$.
  \item $A\subseteq B$: ако $\omega\in A$, то $Z_1(\omega)+
  Z_2(\omega)<c$, т.е.\ $Z_1(\omega) < c-Z_2(\omega)$. Поради
  гъстотата на $\mathbb Q$ в~$\mathbb R$, съществува рационално $q$
  такова, че $Z_1(\omega)<q<c-Z_2(\omega)$; тогава $Z_1(\omega)<q$ и~$Z_2(\omega)<c-q$, т.е.\ $\omega\in B$.
\end{itemize}
Всяко от множествата $\{Z_1<q\}$, $\{Z_2<c-q\}$ е от $\mathcal A$ (по
измеримост на $Z_1,Z_2$), сечението им е от $\mathcal A$, а~изброимото им обединение по $q\in\mathbb Q$ отново е от $\mathcal A$.
Следователно $Z_1+Z_2$ е случайна величина. \qedbg

\textbf{2)} Частен случай на 1) при $b=0$.

\textbf{3)} Записваме произведението алгебрично:
\[
  XY = \tfrac14\big((X+Y)^2-(X-Y)^2\big).
\]
От 1) знаем, че $X\pm Y$ е случайна величина, а~умножението с~константа запазва свойството „случайна величина“ - остава да
докажем, че квадратът на случайна величина също е случайна величина.
Нека $Z$ е произволна случайна величина; показваме $\{Z^2<x\}\in
\mathcal A$ за всяко $x\in\mathbb R$, разглеждайки два случая:
\begin{itemize}[nosep]
  \item $x\le 0$: квадратът на реално число не може да е отрицателен, следователно $\{Z^2<x\}=\varnothing\in\mathcal A$.
  \item $x>0$: множеството $\{Z^2<x\}$ е еквивалентно на $-\sqrt x < Z < \sqrt x$, т.е.
  \[
    \{Z^2<x\} = \{Z<\sqrt x\}\cap\{Z>-\sqrt x\}.
  \]
  Понеже $Z$ е случайна величина, $\{Z<\sqrt x\}\in\mathcal A$; също
  $\{Z>-\sqrt x\}\in\mathcal A$, защото е допълнение на
  $\{Z\le -\sqrt x\}\in\mathcal A$, а~$\sigma$-алгебрите са затворени
  относно допълнение. Сечението на две множества от $\mathcal A$
  също е в~$\mathcal A$, следователно $\{Z^2<x\}\in\mathcal A$.
\end{itemize}
Оттук $XY$, като линейна комбинация на квадрати на случайни
величини, също е случайна величина. \qedbg

\textbf{4)} Записваме частното като произведение: $\dfrac{X}{Y}=X\cdot
\dfrac1Y$. Достатъчно е да покажем, че $\tfrac1Y$ е случайна
величина, т.е.\ $\{\tfrac1Y<x\}\in\mathcal A$ за всяко
$x\in\mathbb R$. Разглеждаме три случая:
\begin{itemize}[nosep]
  \item $x=0$: $\{\tfrac1Y<0\}=\{Y<0\}\in\mathcal A$, защото $Y$ е случайна величина.
  \item $x>0$: $\{\tfrac1Y<x\} = \{Y<0\}\cup\{Y>\tfrac1x\}$ - обединение на две множества от $\mathcal A$, следователно е в~$\mathcal A$.
  \item $x<0$: $\{\tfrac1Y<x\} = \{\tfrac1x<Y<0\} = \{Y>\tfrac1x\}\cap\{Y<0\}$ - сечение на множества от $\mathcal A$, следователно е в~$\mathcal A$.
\end{itemize}
Във всички случаи $\{\tfrac1Y<x\}\in\mathcal A$, следователно
$\tfrac1Y$, а~оттам и~$\tfrac{X}{Y}$, е случайна величина (при
условието $\mathbb P(Y=0)=0$, което гарантира, че делението е
коректно дефинирано с~вероятност~1). \qedbg
\end{proofbox}

\section{Дискретни случайни величини}

\begin{defbox}[~(Индикаторна функция)]
Нека $V=(\Omega,\mathcal A,\mathbb P)$ е вероятностно пространство и~$H\in\mathcal A$. \textbf{Индикаторната функция} $I_H$ се
дефинира като
\[
  I_H(\omega) = \begin{cases} 1, & \omega\in H\\ 0, & \omega\notin H \ (\omega\in H^c)\end{cases},
  \qquad I_H:\Omega\to\{0,1\}.
\]
\end{defbox}

\begin{corbox}[~(свойства на индикаторната функция)]
\begin{enumerate}[label=\arabic*)]
  \item $I_{H^c} = 1-I_H$;
  \item $I_{A\cap B}(\omega) = I_A(\omega)\cdot I_B(\omega)$;
  \item $I_{A\cup B} = I_A+I_B-I_A\cdot I_B$.
\end{enumerate}
\end{corbox}

\begin{thmbox}[]{Лема}
Нека $V=(\Omega,\mathcal A,\mathbb P)$ е вероятностно пространство и~$H\in\mathcal A$. Тогава $I_H$ е случайна величина.
\end{thmbox}

\begin{proofbox}
Полагаме $X(\omega):=I_H(\omega)$ и~разглеждаме изчерпателно
множеството $\{\omega\in\Omega : X(\omega)<c\}$ в~зависимост от
избора на $c\in\mathbb R$ - три случая:
\begin{enumerate}[label=\arabic*)]
  \item $c\le 0$: за $X(\omega)$ възможните стойности са само $0$ и~$1$, така че $X^{-1}((-\infty,c)) = \varnothing \in\mathcal A$ (празното множество е част от всяка $\sigma$-алгебра).
  \item $0<c\le 1$: единственият начин да е изпълнено $X(\omega)<c$ е при $X(\omega)=0$, т.е.\ $X^{-1}((-\infty,c)) = H^c \in\mathcal A$ (защото $H\in\mathcal A$ и~$\sigma$-алгебрите са затворени относно допълнение).
  \item $c>1$: търсим кога функцията е строго под число, което е строго над~$1$. Най-голямата възможна стойност на $X(\omega)$ е $1$, т.е.\ условието е изпълнено винаги: $X^{-1}((-\infty,c)) = \Omega \in\mathcal A$ (по дефиниция на $\sigma$-алгебра). \qedbg
\end{enumerate}
\end{proofbox}

\begin{defbox}[~(Дискретна случайна величина)]
Нека $V=(\Omega,\mathcal A,\mathbb P)$ е вероятностно пространство и~$H=\{H_1,\dots,H_n\}$ (или $\{H_1,H_2,\dots\}$) е \textbf{пълна група
от събития} в~него. Нека са зададени различни реални числа $\bar x =
(x_1,\dots,x_n)$ (или $(x_1,x_2,\dots)$). Тогава
\[
  X(\omega) = X = \sum_j x_j\,I_{H_j}(\omega) = \sum_j x_j\,I_{H_j}
\]
се нарича \textbf{дискретна случайна величина}. От лемата по-горе
всяко $I_{H_j}$ е случайна величина, а~сума от случайни
величини е случайна величина (Свойство~1), следователно $X$ е
случайна величина. Означаваме $H_j=\{X=x_j\}=\{\omega\in\Omega:
X(\omega)=x_j\}$ и~$p_j = \mathbb P(H_j)$, $j\ge 1$.
\end{defbox}

\begin{defbox}[~(Разпределение на дискретна случайна величина)]
Нека $H$ е пълна група от събития във $V$, зададени са $\bar x$ и~$X=\sum_j x_jI_{H_j}$ е дискретна случайна величина. Тогава
таблицата
\[
  \begin{array}{c|cccc}
    X & x_1 & x_2 & \dots & x_j \dots\\\hline
    \mathbb P & p_1 & p_2 & \dots & p_j \dots
  \end{array}
\]
където $\mathbb P(X=x_j)=p_j=\mathbb P(H_j)$ и~$\sum_j p_j=1$, се
нарича \textbf{разпределение} на $X$. Две дискретни случайни
величини имат \emph{еднакво разпределение} (по разпределение,
означено $X\stackrel{d}{=}Y$), ако техните таблици съвпадат, т.е.\
$\mathbb P(X=x_j)=\mathbb P(Y=x_j)$ за всяко $j$.
\end{defbox}

\begin{rembox}[Хистограма]
Графическото представяне на разпределението на дискретна случайна
величина - стълбче с~височина $p_j$ над всяка стойност $x_j$ - се
нарича \textbf{хистограма}.
\end{rembox}

\section{Смяна на променливите}

\begin{thmbox}[]{Теорема}
Нека $X$ е дискретна случайна величина и~$g:\mathbb R\to\mathbb R$.
Тогава $Y=g(X)$ е дискретна случайна величина.
\end{thmbox}

\begin{proofbox}
Ако $X=\sum_i I_{H_i}x_i$, то $Y=\sum_i I_{H_i}
g(x_i)$. Полагайки $y_i:=g(x_i)$, получаваме $Y=\sum_i \mathbb
1_{H_i} y_i$ - дискретна случайна величина по дефиниция. (Ако за
някои $m\neq k$ се окаже $g(x_m)=g(x_k)$, това не е грешка - просто
за удобство обединяваме съответните събития като $H_m\cup H_k$ с~вероятност $p_m+p_k$.) \qedbg
\end{proofbox}

\section{Съвместни дискретни разпределения}

\begin{defbox}[~(Съвместно дискретно разпределение)]
Нека $X,Y$ са две дискретни случайни величини във $V$. Под
\textbf{съвместно дискретно разпределение} на двойката $(X,Y)$
разбираме таблицата
\[
\begin{array}{c|ccccc|c}
 {}_{Y}\diagdown^{X} & x_1 & x_2 & \dots & x_k & \dots & \\\hline
 y_1 & p_{11} & p_{21} & \dots & p_{k1} & \dots & \sum_i p_{i1}\\
 y_2 & p_{12} & p_{22} & \dots & p_{k2} & \dots & \sum_i p_{i2}\\
 \vdots & & & & & & \vdots\\
\end{array}
\]
където $0\le p_{ij}=\mathbb P(X=x_i,\,Y=y_j)$ и~$\sum_{i,j}p_{ij}=1$.
\end{defbox}

\begin{defbox}[~(Маргинални разпределения)]
Разпределенията само на $X$ (респ.\ само на $Y$), получени чрез
сумиране на съвместната таблица по редове/колони -
$\mathbb P(X=x_i)=\sum_j p_{ij}$, $\mathbb P(Y=y_j)=\sum_i p_{ij}$ -
се наричат \textbf{маргинални разпределения}.
\end{defbox}

\begin{defbox}[~(Смяна на променливите на две дискретни случайни величини)]
Нека $X$ и~$Y$ са две дискретни случайни величини във $V$ и~$g:\mathbb R\times\mathbb R\to\mathbb R$. Тогава $Z=g(X,Y)$ се нарича
\textbf{смяна на променливите} на $(X,Y)$ и~е дискретна случайна
величина.
\end{defbox}

\begin{proofbox}
Нека $X=\sum_i I_{H_i}x_i$ по стойности $x_i$ и~$Y=\sum_j I_{\widetilde H_j} y_j$ по стойности $y_j$,
където $H_i=\{X=x_i\}$, $\widetilde H_j = \{Y=y_j\}$. Тогава
\[
  Z = g(X,Y) = \sum_i\sum_j g(x_i,y_j)\,I_{\{X=x_i,Y=y_j\}} = \sum_i\sum_j g(x_i,y_j)\,I_{H_i}I_{\widetilde H_j},
\]
а~групирайки по различните стойности $z=g(x_i,y_j)$, получаваме
$Z=\sum_l z_l\,I_{T_l}$ за подходяща пълна група $\{T_l\}$ -
т.е.\ $Z$ е дискретна случайна величина. \qedbg
\end{proofbox}

\section{Независимост на дискретни случайни величини}

\begin{defbox}[~(Независимост)]
Нека $X,Y$ са дискретни случайни величини във $V$. Казваме, че
$X\perp\!\!\!\perp Y$ (са независими), ако за всички $i,j$:
\[
  \mathbb P(X=x_i,\,Y=y_j) = \mathbb P\big(\{X=x_i\}\cap\{Y=y_j\}\big) \overset{\text{def}}{=} \mathbb P(X=x_i)\cdot\mathbb P(Y=y_j),
\]
или еквивалентно $\mathbb P(X=x_i\mid Y=y_j)=\mathbb P(X=x_i)$ за
всички $i,j$, за които $\mathbb P(Y=y_j)>0$.
\end{defbox}

\begin{defbox}[~(Независимост в~съвкупност)]
Нека $X_1,\dots,X_n$ са $n$ дискретни случайни величини във $V$.
Тогава те са \textbf{независими в~съвкупност}, ако за всеки избор на
стойности $x_1,x_2,\dots,x_n\in\mathbb R$ е изпълнено:
\[
  \mathbb P(X_1=x_1,X_2=x_2,\dots,X_n=x_n) = \prod_{i=1}^{n}\mathbb P(X_i=x_i).
\]
\end{defbox}

\section{Функция на разпределение}

\begin{defbox}[~(Функция на разпределение)]
Нека $X$ е случайна величина във $V$. Тогава
\[
  F_X(x) = \mathbb P(X< x),\qquad \forall x\in[-\infty,+\infty],
\]
се нарича \textbf{функция на разпределение} на случайната величина
$X$. За нея винаги е в~сила $F_X(-\infty)=0$ и~$F_X(+\infty)=1$; за
дискретна случайна величина с~нарастващи стойности $x_1<x_2<\dots$,
$F_X$ е стъпаловидна функция, постоянна между съседните стойности и~скачаща с~$p_j$ при преминаване през~$x_j$.
\end{defbox}
